% Draft 2026-11-03, official setup. \todo{} values come from results/official/analysis.json.
\begin{abstract}
FrontierMemory ranks memory systems for coding and knowledge-work agents by the solve rate of an agent that reads the memory.
We split the path from the corpus to the score into five links: corpus, generation, harness, agent and grader.
We audit each link by running the benchmark's public pipeline with a local open-weights model and a few disclosed serving changes. The no-memory baseline covers all 15 tasks and was specified before any counted trial.
Memory changed the per-task pass rate by $+0.01$ (95\% interval $-0.12$ to $+0.14$, one-sided $p=0.50$; 4 tasks up, 4 down). This model showed no measurable gain, and the interval excludes mean gains above 0.14 per task. Only one task improved in every memory run. In an exploratory placebo arm that replaced the files' text, it passed 0 of 3 attempts.
No corpus run resembles the prompt of 8 of 15 tasks, and the reference solution scores 0 on 4 tasks, as the maintainers document. In the grading code, a failed LLM judge scores the same as a failed agent. A task broke when a dependency released. The same memory's entry file is 21--100\% visible depending on the harness. With a slow model, the network proxy and the agent's own timeouts turn calls into failed attempts.
The leaderboard publishes aggregate scores; a no-memory row and per-task results would let readers separate these effects from memory quality. Its top five scores lie within two task-attempts of 45, and under a binomial model a rerun would put the same entry first only about one time in four; three generations of the same memory system with our model passed from 4 to 8 of 45.
We recommend publishing no-memory rows, per-task results and all attempts, and counting proxy cuts and judge failures as infrastructure failures.
\end{abstract}
